
\chapter{Battery Aging and Asset Preservation Proxy}
\label{ch:aging}

This chapter has been narrowed to match the battery artifacts that are
actually locked in the repository.  The current package does \emph{not}
yet contain a full electrochemical aging simulator or an online
state-of-health-triggered recalibration loop.  What it does contain is a
useful battery-domain proxy surface: (i)~a monotone width-multiplier table
indexed by state of health, and (ii)~an avoided-severity asset-preservation
proxy derived from the locked fault-performance benchmark.

\section{Why Aging Still Matters}

Battery safety and battery aging are linked even when the current code
does not model cell chemistry explicitly.  Hidden SOC violations are
operationally undesirable not only because they break a safety envelope,
but also because repeated excursions toward deep over-charge or
over-discharge are exactly the kind of events that would stress a real
asset.  That is enough to justify an asset-preservation proxy chapter,
provided the claim stays bounded to the evidence that exists today.

\section{Locked Aging Calibration Surface}

The first locked artifact is
\texttt{reports/publication/aging\_calibration\_outputs.csv}, which exports a
simple state-of-health to width-multiplier surface.  The current battery
interpretation is conservative: as usable capacity shrinks, the uncertainty
set should widen rather than remain tuned to a healthier battery.

\begin{table}[ht]
\centering
\caption{Locked uncertainty-width multiplier surface from
  \texttt{aging\_calibration\_outputs.csv}.}
\label{tab:aging-multiplier}
\begin{tabular}{rr}
\toprule
\textbf{SOH} & \textbf{Width multiplier} \\
\midrule
1.00 & 1.000 \\
0.98 & 1.032 \\
0.96 & 1.064 \\
0.94 & 1.096 \\
0.92 & 1.128 \\
0.90 & 1.160 \\
0.85 & 1.240 \\
0.80 & 1.320 \\
0.75 & 1.400 \\
0.70 & 1.480 \\
\bottomrule
\end{tabular}
\end{table}

This table is a design surface rather than a validated closed-loop aging
controller.  It establishes the intended calibration direction
(lower SOH $\rightarrow$ wider intervals), but it should not yet be read as a
full theorem-backed or field-validated law.

\section{Asset-Preservation Proxy Table}

The second locked artifact is
\texttt{reports/aging/asset\_preservation\_proxy\_table.csv}, which is
generated by \texttt{scripts/run\_aging\_proxy\_analysis.py}.  That wrapper
starts from the fault-performance benchmark, measures how much severity the
DC3S stack avoids relative to the deterministic baseline, and maps the
avoided severity into a proxy asset-preservation quantity.  The exact
numbers therefore remain benchmark-derived, not chemistry-derived.

\begin{table}[ht]
\centering
\caption{Representative rows from the locked asset-preservation proxy table.}
\label{tab:asset-preservation-proxy}
\begin{tabular}{lrrr}
\toprule
\textbf{Fault} & \textbf{Severity} & \textbf{Avoided Sev.\ (MWh)} &
  \textbf{Annual fade proxy (\%)} \\
\midrule
dropout & 0.20 & 504.18 & 20.17 \\
dropout & 0.30 & 1912.04 & 76.48 \\
delay\_seconds & 5.0 & 2557.58 & 102.30 \\
spike\_sigma & 3.0 & 232.67 & 9.31 \\
\bottomrule
\end{tabular}
\end{table}

The useful meaning of this table is comparative: larger hidden-violation
severity implies more potential asset stress, and the zero-TSVR DC3S runs
remove that stress channel in the benchmark.  The table is \emph{not} a
warranty-grade life model.

\section{Bounded Claim}

The chapter now makes the narrower claim that the current repository can
honestly support:

\begin{itemize}
  \item the battery evidence package contains an aging-aware \emph{proxy}
    calibration surface,
  \item the benchmark package contains an asset-preservation proxy derived
    from avoided hidden-violation severity, and
  \item both artifacts point in the same operational direction: safer
    runtime behavior should preserve battery health better than a controller
    that allows hidden SOC excursions.
\end{itemize}

\section{What Is Not Yet Implemented}

The current battery repo does \emph{not} yet lock the following items:

\begin{itemize}
  \item a real SOH estimator integrated into the runtime loop,
  \item an online conformal recalibrator triggered automatically by battery
    aging, or
  \item a validated ten-year depreciation model grounded in laboratory or
    field battery-aging measurements.
\end{itemize}

Those remain extension targets rather than locked thesis evidence.

\section{Summary}

The aging story is now aligned with the actual code and artifacts.  The repo
contains a real battery asset-preservation proxy surface, but not a final
electrochemical aging controller.  That is still valuable: it turns aging
from aspirational prose into a measurable witness-domain extension while
keeping the claim appropriately bounded.
